\documentclass[10pt,a4paper]{amsart}
\usepackage{amsmath,amssymb,amsthm,graphicx}
% For commutative diagrams we use Xy-pic
\usepackage[all]{xy}

% We use 2cell for 2-commutative diagrams.
\xyoption{2cell}
% Theorem environments.
%
\theoremstyle{plain}
\newtheorem{theorem}[subsection]{Theorem}
\newtheorem{proposition}[subsection]{Proposition}
\newtheorem{lemma}[subsection]{Lemma}

\theoremstyle{definition}
\newtheorem{definition}[subsection]{Definition}
\newtheorem{example}[subsection]{Example}
\newtheorem{exercise}[subsection]{Exercise}
\newtheorem{situation}[subsection]{Situation}

\theoremstyle{remark}
\newtheorem{remark}[subsection]{Remark}
\newtheorem{remarks}[subsection]{Remarks}

\numberwithin{equation}{subsection}

% Macros
%
\def\lim{\mathop{\mathrm{lim}}\nolimits}
\def\colim{\mathop{\mathrm{colim}}\nolimits}
\def\Spec{\mathop{\mathrm{Spec}}}
\def\Hom{\mathop{\mathrm{Hom}}\nolimits}
\def\Ext{\mathop{\mathrm{Ext}}\nolimits}
\def\SheafHom{\mathop{\mathcal{H}\!\mathit{om}}\nolimits}
\def\SheafExt{\mathop{\mathcal{E}\!\mathit{xt}}\nolimits}
\def\Sch{\mathit{Sch}}
\def\Mor{\mathop{\mathrm{Mor}}\nolimits}
\def\Ob{\mathop{\mathrm{Ob}}\nolimits}
\def\Sh{\mathop{\mathit{Sh}}\nolimits}
\def\NL{\mathop{N\!L}\nolimits}
\def\CH{\mathop{\mathrm{CH}}\nolimits}
\def\proetale{{pro\text{-}\acute{e}tale}}
\def\etale{{\acute{e}tale}}
\def\QCoh{\mathit{QCoh}}
\def\Ker{\mathop{\mathrm{Ker}}}
\def\Im{\mathop{\mathrm{Im}}}
\def\Coker{\mathop{\mathrm{Coker}}}
\def\Coim{\mathop{\mathrm{Coim}}}

% Boxtimes
%
\DeclareMathSymbol{\boxtimes}{\mathbin}{AMSa}{"02}

%
% Macros for moduli stacks/spaces
%
\def\QCohstack{\mathcal{QC}\!\mathit{oh}}
\def\Cohstack{\mathcal{C}\!\mathit{oh}}
\def\Spacesstack{\mathcal{S}\!\mathit{paces}}
\def\Quotfunctor{\mathrm{Quot}}
\def\Hilbfunctor{\mathrm{Hilb}}
\def\Curvesstack{\mathcal{C}\!\mathit{urves}}
\def\Polarizedstack{\mathcal{P}\!\mathit{olarized}}
\def\Complexesstack{\mathcal{C}\!\mathit{omplexes}}
% \Pic is the operator that assigns to X its picard group, usage \Pic(X)
% \Picardstack_{X/B} denotes the Picard stack of X over B
% \Picardfunctor_{X/B} denotes the Picard functor of X over B
\def\Pic{\mathop{\mathrm{Pic}}\nolimits}
\def\Picardstack{\mathcal{P}\!\mathit{ic}}
\def\Picardfunctor{\mathrm{Pic}}
\def\Deformationcategory{\mathcal{D}\!\mathit{ef}}
\usepackage{cleveref}
\crefname{originallemma}{original lemma}{original lemmas}
\Crefname{originallemma}{Original lemma}{Original lemmas}
\makeatletter
\newlabel{definition-envelope}{{0GU5}{source}{}{sourcepointer-1}{}}
\newlabel{definition-envelope@cref}{{[originallemma][0][]0GU5}{[1][0][]source}{}{}{}}
\makeatother
\crefname{originallemma}{original lemma}{original lemmas}
\Crefname{originallemma}{Original lemma}{Original lemmas}
\makeatletter
\newlabel{lemma-base-change-envelope}{{0GU7}{source}{}{sourcepointer-2}{}}
\newlabel{lemma-base-change-envelope@cref}{{[originallemma][0][]0GU7}{[1][0][]source}{}{}{}}
\makeatother
\crefname{originallemma}{original lemma}{original lemmas}
\Crefname{originallemma}{Original lemma}{Original lemmas}
\makeatletter
\newlabel{lemma-cap-cp-bivariant}{{0B7H}{source}{}{sourcepointer-3}{}}
\newlabel{lemma-cap-cp-bivariant@cref}{{[originallemma][0][]0B7H}{[1][0][]source}{}{}{}}
\makeatother
\crefname{originallemma}{original lemma}{original lemmas}
\Crefname{originallemma}{Original lemma}{Original lemmas}
\makeatletter
\newlabel{lemma-composition-envelope}{{0GU6}{source}{}{sourcepointer-4}{}}
\newlabel{lemma-composition-envelope@cref}{{[originallemma][0][]0GU6}{[1][0][]source}{}{}{}}
\makeatother
\crefname{originallemma}{original lemma}{original lemmas}
\Crefname{originallemma}{Original lemma}{Original lemmas}
\makeatletter
\newlabel{lemma-envelope-bivariant}{{0GUC}{source}{}{sourcepointer-5}{}}
\newlabel{lemma-envelope-bivariant@cref}{{[originallemma][0][]0GUC}{[1][0][]source}{}{}{}}
\makeatother
\crefname{originallemma}{original lemma}{original lemmas}
\Crefname{originallemma}{Original lemma}{Original lemmas}
\makeatletter
\newlabel{lemma-pre-derived-chern-class}{{0ESZ}{source}{}{sourcepointer-6}{}}
\newlabel{lemma-pre-derived-chern-class@cref}{{[originallemma][0][]0ESZ}{[1][0][]source}{}{}{}}
\makeatother
\crefname{originallemma}{original lemma}{original lemmas}
\Crefname{originallemma}{Original lemma}{Original lemmas}
\makeatletter
\newlabel{remark-restriction-bivariant}{{0F9Z}{source}{}{sourcepointer-7}{}}
\newlabel{remark-restriction-bivariant@cref}{{[originallemma][0][]0F9Z}{[1][0][]source}{}{}{}}
\makeatother
\crefname{originallemma}{original lemma}{original lemmas}
\Crefname{originallemma}{Original lemma}{Original lemmas}
\makeatletter
\newlabel{situation-setup}{{02QL}{source}{}{sourcepointer-8}{}}
\newlabel{situation-setup@cref}{{[originallemma][0][]02QL}{[1][0][]source}{}{}{}}
\makeatother
\title{Existence and uniqueness of the Chern class, Chern character and Chern polynomials of a perfect object trivialised by an envelope}
\author{The Stacks Project Authors}
\date{stacks-project:chow.tex:local-snapshot-20261012; GNU FDL 1.2 or later}
\begin{document}
\maketitle
\noindent\textit{Editorial note.} All delivered mathematical text is the author's own, copied from the pinned Stacks Project chapter file chow.tex; no statement or proof step was rewritten, reordered or removed.\par\medskip
\noindent\textit{Editorial note.} Preamble provenance: the pinned chapter file opens with \textbackslash{}input\{preamble\} and carries no preamble text of its own, so the project's single Stacks-wide preamble staged at sources/stacks-chow/preamble.tex is excerpted verbatim as two spans, lines 14-18 (Xy-pic setup) and lines 154-217 (the native theorem declarations and mathematical macros). Nothing was added, removed or redefined; the wrapper class and the amsmath/amssymb/amsthm/graphicx base packages are supplied by the assembler.\par\medskip
\noindent\textit{Editorial note.} Same-chapter pointers are located at their \textbackslash{}label lines inside the pinned chapter file; cross-chapter pointers are located in the project's own tag list staged in the snapshot (sources/stacks-chow/tags.txt), at the verbatim lines carrying their canonical tags.\par\medskip
\noindent\textit{Editorial note.} The author's opening sentence at source lines 9230-9231, 'We will prove the lemma for the chern class c\_p(E) in A\textasciicircum{}p(X) and omit the arguments for the other cases.', is retained verbatim at the start of the counted proof. The construction and both uniqueness arguments are given in full for the Chern class; the Chern character and Chern polynomials are the same argument with the corresponding additivity lemma, as the statement announces.\par\medskip
\noindent\textit{Allowed standard prerequisites (editorial).}\begin{itemize}
\item Independence of the Chern class of a bounded finite locally free complex representing a perfect object. Same chapter chow.tex, label lemma-pre-derived-chern-class, invoked at proof line 9239. Makes the class c\_p of a complex representing Lg\textasciicircum{}*E well defined on any locally finite type T -> X, which is what allows the construction to be base changed freely.
\item Restriction of bivariant classes and the bivariant character of c\_p. Same chapter chow.tex, labels remark-restriction-bivariant (proof line 9246) and lemma-cap-cp-bivariant (proof line 9248). Identifies the class obtained after a further base change with the restriction of the earlier one, so that the two restrictions in the envelope descent lemma really coincide.
\item Descent of a bivariant class along an envelope, and the envelope properties of base change and composition. Same chapter chow.tex, labels lemma-envelope-bivariant (proof lines 9253 and 9268), lemma-base-change-envelope and lemma-composition-envelope (proof lines 9266-9267), and definition-envelope (statement, proof line 9217). Produces the descended bivariant class c\_p(E) on X from the class on Y, and gives the uniqueness by applying the same descent lemma to the product envelope Y times\_X Y' -> X.
\end{itemize}\medskip
\ifdefined\hypertarget\hypertarget{sourcepointer-1}{}\fi
\noindent\textit{Source reference.} Original lemma 0GU5: sources/stacks-chow/pinned-original-source.tex, source lines 3553--3553. Navigation pointer for \textbackslash{}ref\{definition-envelope\} used inside the counted span (line 9217); the argument that consumes it is retained in the span.\par\medskip
\ifdefined\hypertarget\hypertarget{sourcepointer-2}{}\fi
\noindent\textit{Source reference.} Original lemma 0GU7: sources/stacks-chow/pinned-original-source.tex, source lines 3580--3580. Navigation pointer for \textbackslash{}ref\{lemma-base-change-envelope\} used inside the counted span (line 9266); the argument that consumes it is retained in the span.\par\medskip
\ifdefined\hypertarget\hypertarget{sourcepointer-3}{}\fi
\noindent\textit{Source reference.} Original lemma 0B7H: sources/stacks-chow/pinned-original-source.tex, source lines 7533--7533. Navigation pointer for \textbackslash{}ref\{lemma-cap-cp-bivariant\} used inside the counted span (line 9248); the argument that consumes it is retained in the span.\par\medskip
\ifdefined\hypertarget\hypertarget{sourcepointer-4}{}\fi
\noindent\textit{Source reference.} Original lemma 0GU6: sources/stacks-chow/pinned-original-source.tex, source lines 3567--3567. Navigation pointer for \textbackslash{}ref\{lemma-composition-envelope\} used inside the counted span (line 9267); the argument that consumes it is retained in the span.\par\medskip
\ifdefined\hypertarget\hypertarget{sourcepointer-5}{}\fi
\noindent\textit{Source reference.} Original lemma 0GUC: sources/stacks-chow/pinned-original-source.tex, source lines 6707--6707. Navigation pointer for \textbackslash{}ref\{lemma-envelope-bivariant\} used inside the counted span (line 9253); the argument that consumes it is retained in the span.\par\medskip
\ifdefined\hypertarget\hypertarget{sourcepointer-6}{}\fi
\noindent\textit{Source reference.} Original lemma 0ESZ: sources/stacks-chow/pinned-original-source.tex, source lines 9091--9091. Navigation pointer for \textbackslash{}ref\{lemma-pre-derived-chern-class\} used inside the counted span (line 9239); the argument that consumes it is retained in the span.\par\medskip
\ifdefined\hypertarget\hypertarget{sourcepointer-7}{}\fi
\noindent\textit{Source reference.} Original lemma 0F9Z: sources/stacks-chow/pinned-original-source.tex, source lines 6184--6184. Navigation pointer for \textbackslash{}ref\{remark-restriction-bivariant\} used inside the counted span (line 9246); the argument that consumes it is retained in the span.\par\medskip
\ifdefined\hypertarget\hypertarget{sourcepointer-8}{}\fi
\noindent\textit{Source reference.} Original lemma 02QL: sources/stacks-chow/pinned-original-source.tex, source lines 1372--1372. Navigation pointer for \textbackslash{}ref\{situation-setup\} used inside the counted span (line 9214); the argument that consumes it is retained in the span.\par\medskip
\begin{lemma}
\label{lemma-defined-by-envelope}
Let $(S, \delta)$ be as in Situation \ref{situation-setup}.
Let $X$ be locally of finite type over $S$. Let $E \in D(\mathcal{O}_X)$
be a perfect object. Assume there exists an envelope
$f : Y \to X$ (Definition \ref{definition-envelope})
such that $Lf^*E$ is isomorphic in $D(\mathcal{O}_Y)$
to a locally bounded complex $\mathcal{E}^\bullet$ of finite locally free
$\mathcal{O}_Y$-modules. Then there exists unique bivariant classes
$c(E) \in \prod_{p \geq 0} A^p(X)$,
$ch(E) \in \prod_{p \geq 0} A^p(X) \otimes \mathbf{Q}$, and
$P_p(E) \in A^p(X)$, independent of the choice of $f : Y \to X$
and $\mathcal{E}^\bullet$, such that the restriction of these classes
to $Y$ are equal to $c(\mathcal{E}^\bullet)$,
$ch(\mathcal{E}^\bullet)$, and $P_p(\mathcal{E}^\bullet)$.
\end{lemma}

\begin{proof}
Fix $p \in \mathbf{Z}$. We will prove the lemma for the chern
class $c_p(E) \in A^p(X)$ and omit the arguments for the other cases.

\medskip\noindent
Let $g : T \to X$ be a morphism locally of finite type such that
there exists a locally bounded complex $\mathcal{E}^\bullet$ of finite locally
free $\mathcal{O}_T$-modules representing $Lg^*E$ in $D(\mathcal{O}_T)$.
The bivariant class $c_p(\mathcal{E}^\bullet) \in A^p(T)$
is independent of the choice of $\mathcal{E}^\bullet$ by
Lemma \ref{lemma-pre-derived-chern-class}.
Let us write $c_p(Lg^*E) \in A^p(T)$ for this class.
For any further morphism $h : T' \to T$ which is locally
of finite type, setting $g' = g \circ h$ we see that
$L(g')^*E = L(g \circ h)^*E = Lh^*Lg^*E$ is represented by
$h^*\mathcal{E}^\bullet$ in $D(\mathcal{O}_{T'})$.
We conclude that $c_p(L(g')^*E)$ makes sense and is equal to the
restriction (Remark \ref{remark-restriction-bivariant})
of $c_p(Lg^*E)$ to $T'$ (strictly speaking this requires an application of
Lemma \ref{lemma-cap-cp-bivariant}).

\medskip\noindent
Let $f : Y \to X$ and $\mathcal{E}^\bullet$ be as in the statement
of the lemma. We obtain a bivariant class $c_p(E) \in A^p(X)$ from
an application of Lemma \ref{lemma-envelope-bivariant}
to $f : Y \to X$ and the class $c' = c_p(Lf^*E)$ we constructed
in the previous paragraph. The assumption in the lemma is
satisfied because by the discussion in the previous paragraph we have
$res_1(c') = c_p(Lg^*E) = res_2(c')$ where
$g = f \circ p = f \circ q : Y \times_X Y \to X$.

\medskip\noindent
Finally, suppose that $f' : Y' \to X$ is a second envelope
such that $L(f')^*E$ is represented by a bounded complex of
finite locally free $\mathcal{O}_{Y'}$-modules. Then it follows
that the restrictions of $c_p(Lf^*E)$ and $c_p(L(f')^*E)$
to $Y \times_X Y'$ are equal. Since $Y \times_X Y' \to X$
is an envelope (Lemmas \ref{lemma-base-change-envelope} and
\ref{lemma-composition-envelope}), we see that our two candidates for $c_p(E)$
agree by the unicity in Lemma \ref{lemma-envelope-bivariant}.
\end{proof}

\end{document}
